\section{结论}
\label{sec:conclusion}

从点光源视角光栅化高斯，可以记录每个图元截获的通量比例。
\ours 将这一沉积量转化为可学习图集，只需为每个光源额外执行一次光栅化，即可为任意接收点提供经过滤波、可进一步修正的阴影检验和光传输项。
结合按辐射度量域分阶段的优化，该方法在三个基准上表现最为稳定一致，训练时间仅为常规方法的一小部分；其相对于光源定义的设计能否在未见过的光照方向上带来优势，仍有待计划中的实验回答。
